% Глава: 09-vim
% Статус: наполнено (первый прогон).
\chapter{Текстовый редактор Vim}
\label{ch:09-vim}

\section{Идея и моды}

\cmd{vim} (vi editor) --- клавиатурный редактор, установленный
практически в любой Linux-системе. Ключевая идея: \textbf{моды}.
Не всё время вы в режиме «печатаю текст»: есть отдельные режимы
для навигации, команд и ввода.

Vim работает в терминале вашей ВМ на \fpath{pglab.nsu.ru} --- это
основной способ править файлы на сервере без графического редактора.

\begin{center}
\begin{tabular}{lp{0.7\linewidth}}
\toprule
\textbf{Режим} & \textbf{Назначение} \\
\midrule
Normal   & навигация и вызов команд; режим по умолчанию \\
Insert   & ввод текста (аналог «обычного» редактора) \\
Visual   & выделение фрагмента для операций \\
Command  & строка команд в самом низу (\cmd{:}) для записи, выхода и т.п. \\
\bottomrule
\end{tabular}
\end{center}

Наиболее частая ошибка новичка --- попробовать печатать в режиме
Normal: текст не появляется. Чтобы начать печатать, нужно
войти в Insert (\cmd{i} --- вставить перед курсором, \cmd{a} ---
после, \cmd{o} --- новая строка снизу). Вернуться в Normal ---
нажать \cmd{Esc}.

\section{Базовая навигация (Normal mode)}

Перемещение без мыши:
\begin{minted}[bgcolor=gray!6]{bash}
h  j  k  l      # влево, вниз, вверх, вправо
w  b            # следующее / предыдущее слово
0  $            # начало / конец строки
gg  G           # начало / конец файла
Ctrl-d Ctrl-u   # полстраницы вниз / вверх
\end{minted}

Поиск и переход:
\begin{minted}[bgcolor=gray!6]{bash}
/pattern        # поиск вперёд от курсора
n  N            # следующее / предыдущее совпадение
\end{minted}

\section{Операции над текстом (Normal mode)}

Vim строит операции как \emph{глагол + объект}: \cmd{d} (delete),
\cmd{y} (yank/copy), \cmd{p} (paste), \cmd{c} (change) плюс
объект: \cmd{w} (word), \cmd{d} (line), \cmd{i}+пары (inside
parentheses/quotes), \cmd{t} (tag) и т.д.
\begin{minted}[bgcolor=gray!6]{bash}
dd         # удалить текущую строку
yy         # скопировать текущую строку
p          # вставить после курсора
u          # отменить
.          # повторить последнюю команду
dw         # удалить слово
ci"        # изменить внутри кавычек
\end{minted}

Операции можно умножать: \cmd{3dd} --- удалить 3 строки,
\cmd{5w} --- вперёд 5 слов.

\section{Запись и выход (Command mode)}

Внизу нажимаете \cmd{:} и вводите команду:
\begin{minted}[bgcolor=gray!6]{bash}
:w          # сохранить
:w new.txt  # сохранить как new.txt
:wq         # сохранить и выйти
:q          # выйти; не выйдет при несохранённых изменениях
:q!         # выйти, отбросив изменения
\end{minted}

Типовой сценарий: \cmd{Esc} для возврата в Normal, затем
\cmd{:wq} --- сохранить и выйти.

\section{Начальная настройка}

Vim читает \fpath{~/.vimrc} --- ваш конфиг.
Минимально полезный набор:
\begin{minted}[bgcolor=gray!6]{bash}
" show line numbers
set number
" 4 spaces, no tabs
set expandtab tabstop=4 shiftwidth=4
" highlight the current line
set cursorline
" sensible defaults for a modern system
set nocompatible
\end{minted}

Для работы с проектами, синтаксисом и автодополнением обычно
используют более развитые среды (Neovim, Emacs, IDE) --- но
навык базового Vim остаётся полезным на удалённых серверах,
в контейнерах и в rescue-режимах, где тяжёлый редактор
не поставить.

\section{Типичные ошибки}
\begin{itemize}
  \item Текст не печатается --- вы в Normal; нажмите \cmd{i}.
  \item \cmd{:q} не выходит --- есть несохранённые изменения;
    используйте \cmd{:q!} или сначала \cmd{:w}.
  \item Удалили случайно много строк --- \cmd{u} (undo) или
    \cmd{u} несколько раз; \cmd{U} --- восстановить строку целиком.
  \item Перепутали \cmd{d} (delete) и \cmd{y} (yank) ---
    данные потеряны? \cmd{p} вставит то, что было в буфере.
  \item Не нашли нужную строку --- \cmd{/pattern}, затем \cmd{n}.
\end{itemize}

\section{Связь с остальным курсом}

Vim удобен для правки конфигов
(\fpath{/etc/hosts}, \fpath{.bashrc}), Makefile
(\nameref{ch:14-make}), Dockerfile
(\nameref{ch:20-dockerfile}), C-кода
(\nameref{ch:13-c-compilation}). Внутри tmux
(\nameref{ch:10-tmux}) Vim --- основной редактор;
внутри контейнеров (\nameref{ch:19-docker-basics}) он часто
единственный доступный.

\section{Задачи для самопроверки}
\begin{enumerate}
  \item Объясните, чем отличаются \cmd{i}, \cmd{a}, \cmd{o}
    для входа в Insert.
  \item Как удалить текущую строку и как её восстановить?
  \item Как найти все вхождения слова и переключаться между
    ними?
  \item Чем \cmd{:q} отличается от \cmd{:q!}?
  \item Как в Vim перейти в конец файла, затем удалить
    последние три строки?
\end{enumerate}
